% ======================================================================
% Experiments (restructured 2026-09-26 around RQ1-RQ3; approved plan).
% Tables 2 and 3 are native LaTeX tables generated from
% paper/trajqa_b_results_20260925_1300/{results_table.md, figure/Ablation.md}
% (identical to the package's PNG tables, cross-checked against the CSVs);
% Table 4 is parsed from TRAJQA_B_SPEC.md §21.8 (held-out block).
% Protocol details, full tables and secondary diagnostics: Appendix app:exp_setup.
% ======================================================================
\section{Experiments}
\label{sec:experiments}

\paragraph{Dataset and evaluation.}
\label{sec:qa}
TrajQA-B draws on 42,151 curated Stereo4D scenes and uses source-video-disjoint QA splits: 56,120 training, 1,318 validation, 4,492 test, and 3,430 additional held-out questions (Table~\ref{tab:qa-splits}). Seven question families cover camera and object displacement, endpoint displacement per elapsed second, interval comparisons, motion source, and duration (Appendix~\ref{app:qa}). The three four-choice numeric families contribute 7,500 evaluation questions and have 25.0\% uniform-guessing accuracy. In every table \emph{speed} means endpoint displacement divided by elapsed time, rather than path length per time. Gate~0 retains questions whose reference answer a deterministic reader can reproduce from the capped records; this tests record reading and calculation, not reconstruction accuracy. Option ranks and binary answers are balanced, then audited against question-only heuristics. The separately selected stage-1 caption training scenes also exclude QA evaluation scenes; its source-video overlap has not been audited (Appendix~\ref{app:training}). Appendix~\ref{app:exp_setup} gives the full protocols and denominators.

\paragraph{Implementation details.}
Our backbone is Qwen3-VL-8B-Instruct with frozen vision and base language weights. The trainable trajectory projector has 85.93M parameters; language-model LoRA adapters have rank 128, scale 256, dropout 0.05, and 349.18M parameters. Each sample contains eight uniformly sampled video frames, resized to a longest side of 384 pixels (144 visual tokens per frame), the object-labelled trajectory stream capped at 512 records, and a question or instruction. We use three stages: 26,000-step projector warm-up on 64,523 rule-generated motion captions from a separately selected caption corpus; 8,400-step projector-plus-LoRA QA training with trajectory tokens but no frames; and 3,507-step joint QA training with frames restored and independently dropped with probability 0.2. All stages use AdamW with $\beta=(0.9,0.95)$, zero weight decay, gradient clipping at 1.0, global batch size 16, and seed 0. The peak projector learning rate is $2\times10^{-4}$ throughout; the adapters use $10^{-4}$ in stages 2 and 3. We train on the assistant's evidence and answer text with next-token cross-entropy and decode greedily. The frames-only Qwen3-VL-8B control was fine-tuned on the same QA training split, although its exact schedule and checkpoint selection are unavailable in the manuscript. Most external baselines are zero-shot and some use their official video pipelines or a smaller proprietary-model subset; they contextualize performance rather than isolate the effect of trajectory input. Stage-1 captions use metadata frame rate while QA targets use frame timestamps, a source of noisy time labels described in Appendix~\ref{app:data_curation}.

% ----------------------------------------------------------------------
\subsection{Comparisons on TrajQA-B}
\label{sec:exp-main}

\begin{table}[t]
\caption{TrajQA-B numeric accuracy (\%). Numeric-3 pools the three four-choice families ($n=7{,}500$; guessing: 25.0); split and family scores are also shown. External models are zero-shot except MLLM-4D-SFT, while ours is fine-tuned on the TrajQA-B training split. Proprietary models use a 1,000-question subset. Bold marks the largest displayed score in each column; protocols and subsets differ. Full results appear in Appendix~\ref{app:exp_setup}.}
\label{tab:main}
\begin{center}
\small
\setlength{\tabcolsep}{3.6pt}
\begin{tabular}{lcccccc}
\toprule
\textbf{Model} & \textbf{Numeric-3}$\uparrow$ & \textbf{Held-out}$\uparrow$ & \textbf{Test}$\uparrow$ & \textbf{Cam.\ disp.}$\uparrow$ & \textbf{Obj.\ disp.}$\uparrow$ & \textbf{Rate}$\uparrow$ \\
\midrule
\rowcolor{blue!12}\multicolumn{7}{l}{\emph{Proprietary models (1,000-question subset)}} \\
GPT-6 Astra & 39.2 & 37.9 & 40.5 & 46.0 & 32.7 & 33.9 \\
Claude Fable 5.1 & 35.2 & 35.8 & 34.5 & 39.6 & 35.2 & 26.9 \\
\midrule
\rowcolor{teal!12}\multicolumn{7}{l}{\emph{Open-source base models, 7B and 8B}} \\
Qwen2.5-VL-7B & 27.3 & 26.6 & 27.9 & 27.1 & 27.7 & 27.3 \\
Qwen3-VL-8B & 28.1 & 28.5 & 27.7 & 28.9 & 27.9 & 26.9 \\
LLaVA-NeXT-Video-7B & 25.0 & 24.5 & 25.4 & 25.7 & 24.2 & 24.7 \\
InternVideo2.5-8B & 25.5 & 25.5 & 25.5 & 25.5 & 24.9 & 26.1 \\
InternVL2.5-8B & 27.0 & 27.3 & 26.8 & 27.7 & 27.4 & 25.5 \\
InternVL3.5-8B & 26.8 & 25.8 & 27.5 & 26.3 & 26.4 & 27.9 \\
\midrule
\rowcolor{teal!12}\multicolumn{7}{l}{\emph{Open-source base models, 32B}} \\
Qwen2.5-VL-32B & 28.4 & 28.3 & 28.5 & 29.7 & 27.3 & 27.5 \\
Qwen3-VL-32B & 27.8 & 27.8 & 27.7 & 27.9 & 25.9 & 29.5 \\
\midrule
\rowcolor{teal!12}\multicolumn{7}{l}{\emph{Open-source base models, 38B}} \\
InternVL2.5-38B & 29.4 & 29.3 & 29.5 & 30.3 & 29.9 & 27.5 \\
InternVL3.5-38B & 28.5 & 28.9 & 28.3 & 29.4 & 27.5 & 28.2 \\
\midrule
\rowcolor{violet!12}\multicolumn{7}{l}{\emph{3D spatial-reasoning models}} \\
VG-LLM (Qwen2.5-VL-7B) & 25.2 & 24.6 & 25.7 & 24.8 & 25.4 & 25.6 \\
VLM-3R (LLaVA-Video-7B) & 26.9 & 26.9 & 27.0 & 27.1 & 27.0 & 26.7 \\
\midrule
\rowcolor{orange!12}\multicolumn{7}{l}{\emph{4D instruction-tuned model (Stereo4D source)}} \\
MLLM-4D-SFT (Qwen3-VL-8B) & 25.0 & 24.7 & 25.3 & 25.3 & 24.9 & 24.8 \\
\midrule
\rowcolor{green!15}\multicolumn{7}{l}{\emph{TrajQA-B fine-tuned model}} \\
Ours (frames + trajectory tokens) & \textbf{82.2} & \textbf{81.7} & \textbf{82.5} & \textbf{89.4} & \textbf{77.6} & \textbf{75.2} \\
\bottomrule
\end{tabular}
\end{center}
\end{table}

Table~\ref{tab:main} places our model next to 15 external baselines. Open-source and specialised 3D models score between 25.0\% and 29.4\% on the pooled numeric questions; the two proprietary models reach 39.2\% and 35.2\% on their smaller subset. Comparisons from 7B to 38B do not show a clear gain from parameter count in this collection of models, although architectures and video protocols also vary.
MLLM-4D-SFT scores 25.0\% after training on a separate 4D QA corpus drawn from the same source videos; 85\% of evaluation clips appear in its SFT data. Its strong option-position preference provides a plausible explanation for the poor transfer to our balanced options (Appendix~\ref{app:exp_setup}).
Our fine-tuned model reaches 82.2\% overall, with 81.7\% on held-out and 82.5\% on test. A frames-only Qwen3-VL-8B arm trained on the same question split reaches 48.3\% (Table~\ref{tab:ablation}). This task-adapted visual reference lacks timestamps and may omit the queried endpoint frames; its undocumented training schedule also prevents a controlled modality comparison.

% ----------------------------------------------------------------------
\subsection{Input Ablation Studies}
\label{sec:exp-ablation}

\begin{table}[t]
\caption{Input interventions on the full evaluation set. All rows except frames only use the same final model; that row uses a separately fine-tuned visual model. Bold marks the best value in each column, including ties. Donor tokens replace trajectory content with that of another scene while preserving format, count, and timestamps. Cam./Obj.: camera/object displacement; rate: endpoint displacement per elapsed second. All results use one training seed.}
\label{tab:ablation}
\begin{center}
\footnotesize
\setlength{\tabcolsep}{3pt}
\begin{tabular}{lcccccccc}
\toprule
\textbf{Inputs at inference} & \textbf{Frames} & \textbf{Tokens} & \textbf{Numeric-3}$\uparrow$ & \textbf{Held-out}$\uparrow$ & \textbf{Test}$\uparrow$ & \textbf{Cam.}$\uparrow$ & \textbf{Obj.}$\uparrow$ & \textbf{Rate}$\uparrow$ \\
\midrule
Question only & No & No & 26.8 & 27.5 & 26.2 & 27.2 & 27.0 & 25.9 \\
Frames only (separate FT) & \checkmark & No & 48.3 & 47.8 & 48.8 & 53.7 & 44.5 & 43.9 \\
Frames + donor tokens & \checkmark & donor & 29.6 & 30.6 & 28.8 & 30.7 & 28.8 & 28.6 \\
Trajectory tokens only & No & \checkmark & \textbf{82.2} & 81.5 & \textbf{82.8} & \textbf{89.4} & \textbf{77.7} & \textbf{75.5} \\
Frames + traj.\ tokens (full) & \checkmark & \checkmark & \textbf{82.2} & \textbf{81.7} & 82.5 & \textbf{89.4} & 77.6 & 75.2 \\
\bottomrule
\end{tabular}
\end{center}
\end{table}

\begin{table}[t]
\caption{Metric grounding of generated evidence on held-out numeric questions ($n=3{,}312$). Acc.: accuracy; $E_{\mathrm{obj}}$ / $E_{\mathrm{cam}}$: median position error. Bold marks the best value for each metric, including ties. Within $1.2\times$: fraction whose stated camera displacement / object displacement / endpoint rate differs from the reference by less than a multiplicative factor of 1.2.}
\label{tab:evidence}
\begin{center}
\footnotesize
\setlength{\tabcolsep}{3pt}
\begin{tabular}{lcccc}
\toprule
\textbf{Inputs at inference} & \textbf{Acc.}$\uparrow$ & $\boldsymbol{E_{\mathrm{obj}}}$ \textbf{(m)}$\downarrow$ & $\boldsymbol{E_{\mathrm{cam}}}$ \textbf{(m)}$\downarrow$ & \textbf{Within $1.2\times$ $\uparrow$: cam / obj / rate} \\
\midrule
Question only & 27.5 & 3.56 & 0.83 & 21.0 / 15.5 / 14.9 \\
Frames + donor tokens & 30.6 & 3.44 & 1.35 & 18.2 / 17.7 / 16.0 \\
Trajectory tokens only & 81.5 & \textbf{0.32} & \textbf{0.14} & \textbf{91.4} / \textbf{80.0} / \textbf{80.3} \\
Frames + trajectory tokens (full) & \textbf{81.7} & 0.33 & \textbf{0.14} & 91.1 / \textbf{80.0} / 79.9 \\
\bottomrule
\end{tabular}
\end{center}
\end{table}

\textbf{Trajectory content changes the answers.}
The same final model scores 26.8\% with the question alone and 29.6\% with frames plus another scene's tokens, versus 82.2\% with the correct tokens (Table~\ref{tab:ablation}). Both controls remain above uniform guessing, so they do not establish the complete absence of shortcuts; their large gap from the correct-token condition shows that trajectory content matters. The separately trained frames-only arm reaches 48.3\%, roughly 34 points below the full model. On held-out questions, the donor-token intervention lowers accuracy from 81.7\% to 30.6\%. The final model without frames also scores 82.2\% (Appendix~\ref{app:exp_setup}), indicating that this benchmark mainly probes use of supplied metric tracks rather than added value from fusing frames with them.

\textbf{Generated evidence tracks metric quantities.}
On held-out questions, the positions stated in the model's evidence differ from the reference by a median of 0.33~m for objects and 0.14~m for the camera (Table~\ref{tab:evidence}). Its stated displacement or speed lies within a factor of 1.2 of the reference on 80\% to 91\% of questions. With donor tokens, median position errors rise to metres and these rates fall to 16\% to 18\%. The evidence analysis supports use of the trajectory values beyond selecting an option, subject to the reference-track quality and parsing protocol in Appendix~\ref{app:exp_setup}. Qualitative examples appear in Appendix~\ref{app:qualitative}.

\textbf{Module-facing diagnostics.}
Selected scenes average 44.8 keypoints; the tokenizer uses a median of 165 records and its intermediate spline fits meet the target on 99.82\% of pieces. Donor substitution probes use of correct records. All three main numeric families can be answered from endpoint positions and times, so these diagnostics do not establish a benefit from part labels, event splitting, adaptive fitting, or projector warm-up.

\textbf{Where the improvement applies.}
Camera displacement is the easiest numeric family for the full model (89.4\%), followed by object displacement (77.6\%) and endpoint rate (75.2\%). The final model has the same pooled score with tokens alone as with tokens and frames (82.2\%), so these results demonstrate use of reconstructed motion, not added value from visual fusion. The 81.7\% held-out and 82.5\% test scores show similar performance across the two QA splits, but stage-1 source-video overlap has not been audited. Because donor input remains above 25\% chance and the visual arm lacks queried timestamps, neither control rules out all shortcuts or supplies a fully matched visual comparison.

\paragraph{Evaluation composition.}
Table~\ref{tab:qa-splits} shows the family denominators and the number of scenes and source videos in each split. The held-out and test splits contain 3,292 camera-displacement, 2,172 object-displacement, and 2,036 endpoint-rate questions; these are the 7,500 questions underlying the pooled numeric score. The other four families contribute 422 evaluation questions in total. Their reported aggregate in Appendix~\ref{app:exp_setup} is therefore a smaller, heterogeneous diagnostic, not evidence that the proposed event and part representations improve interior-motion reasoning.

\begin{table}[H]
\caption{TrajQA-B split counts (descriptive, not performance metrics). Splits are disjoint in source videos. Cam./Obj./Rate count the three 4-way families (chance 25\%); Other pools interval\_disp, interval\_speed, motion\_source, and binary motion\_duration. Rate denotes endpoint displacement divided by elapsed time.}
\label{tab:qa-splits}
\begin{center}
\footnotesize
\setlength{\tabcolsep}{3pt}
\begin{tabular}{lrrrrrrr}
\toprule
\textbf{Split} & \textbf{Questions} & \textbf{Cam.} & \textbf{Obj.} & \textbf{Rate} & \textbf{Other} & \textbf{Scenes} & \textbf{Videos} \\
\midrule
train & 56,120 & 21,240 & 16,492 & 14,980 & 3,408 & 12,409 & 2,198 \\
validation & 1,318 & 524 & 404 & 340 & 50 & 668 & 162 \\
dev test & 4,492 & 1,692 & 1,284 & 1,212 & 304 & 2,241 & 578 \\
held-out test & 3,430 & 1,600 & 888 & 824 & 118 & 862 & 166 \\
\bottomrule
\end{tabular}
\end{center}
\end{table}

\noindent\textbf{Residual data risks.}
Because 98.4\% of clips share a source video, the unaudited stage-1 source-video intersection remains consequential. Caption frame rates differ from QA timestamps by more than 10\% in 12.85\% of scenes, potentially adding rate-label noise.
